\begin{comment}
\begin{abstract}
    In this paper, we prove the existence of a finite step multi-round language model set $\mathrm{C} = \{\Phi_i = \phi_{i_1}\bullet \phi_{i_2}\bullet\cdots\bullet \phi_{i_m} :\mathbb{R}^d \to \mathbb{R}^d | m \leq d^2\}$ for universal approximation. That is, for any given continuous mapping $f = f_m\circ \cdots \circ f_1: \mathbb R^d \to \mathbb R^d$, compact domain $\Omega$ and $\varepsilon$, there is a multi-round sentences that $\Phi \in C$ that approximate $f$ on $\Omega$ with error less than $\varepsilon$. Furthermore, we prove that a regularization on middle step so that $\phi_{i_k}$ approximate $f_k$ will benefit the overall approximation. We show that any Turing computable function can be approximated in finite rounds of language model interaction.

\end{abstract}
\end{comment}

\begin{abstract}
 Recent advancements in cognitive science and multi-round reasoning techniques for Large Language Models (LLMs) suggest that iterative thinking processes improve problem-solving performance in complex tasks. Inspired by this, approaches like Chain-of-Thought, debating, and self-refinement have been applied to auto-regressive LLMs, achieving significant successes in tasks such as mathematical reasoning, commonsense reasoning, and multi-hop question answering. Despite these successes, the theoretical basis for how multi-round reasoning enhances problem-solving abilities remains underexplored.
In this work, we investigate the approximation, learnability, and generalization properties of multi-round auto-regressive models. We show that Transformers with finite context windows are universal approximators for steps of Turing-computable functions and can approximate any Turing-computable sequence-to-sequence function through multi-round reasoning. We extend PAC learning to sequence generation and demonstrate that multi-round generation is learnable even when the sequence length exceeds the model's context window. 
Finally, we examine how generalization error propagates across rounds, and show how the aforementioned approaches can help constrain this error, ensuring outputs stay within an expectation boundary. This work sheds light on the systemic theoretical foundations of multi-round sequence learning and reasoning, emphasizing its role in inference complexity.
\end{abstract}